\chapter*{Conferencia XVI}
\addcontentsline{toc}{chapter}{LECTURA I - El timbre y el carácter del violín}
SEÑORES: En esta hora se presenta el tema de la máxima potencia tonal del violín. Hoy el
La mayor potencia del sonido del violín atrae más atención de usuarios, constructores y estudiantes de violín que en cualquier otro período de la historia del violín. La causa de este hecho se encuentra en la mayor capacidad de los auditorios modernos. El hombre es un animal gregario y amante de la música; y, por alguna razón oculta, su disfrute de la música es como el cuadrado de su número. Sin duda, este hecho se debe a la progresión geométrica de la acción mental simpática. Así, cuando solo hay un oyente presente en el ensayo, la acción simpática permanece en cero; un hecho demostrado por la notable ausencia de demostración. Al agregarse otro oyente, la acción simpática se desarrolla como el cuadrado de dos; con lo cual, el oyente único se vuelve consciente de un disfrute adicional. Para el músico, tal aumento en el disfrute por parte del oyente es un estímulo mayor que el
estímulo del oro fino.
Es siempre un objetivo de la ambición humana lograr grandes hazañas. Es la ambición del violinista solista complacer a sus clientes en los auditorios más grandes por igual que a los clientes en los auditorios más pequeños. Tal ambición es loable, pero las dificultades se acercan rápidamente a lo insuperable. De ahí el "alto" precio de los violines de gran sonido. Pero la demanda actual de violines de "gran sonido" es una debilidad innecesaria en los usuarios de violín. De los miles que dedican su máxima energía al violín,
235

Pero solo un número limitado será llamado a presentarse en los auditorios más grandes. Solo el genio, combinado con una dedicación casi sobrehumana, puede garantizarle al violinista solista de hoy el mayor reconocimiento. Por lo tanto, las oportunidades para violines de sonido potente son limitadas; pero es nuestra debilidad admirar a quien dice: "Tu violín tiene un sonido potente".
Pis dijo: «Debemos aceptar a la humanidad tal como la encontramos»; pero debido al permiso implícito para no dejarla como la encontramos, por lo tanto, afirmo que es un halago de mucho mayor valor cuando una persona competente dice: «Tu violín tiene un sonido hermoso». Tal afirmación se basa en el hecho de que la música es, ante todo, bella. No quiero decir que un sonido de violín "grande" no pueda ser hermoso; sino que la distancia del oyente debe ser proporcional a la "grandeza" del sonido. En auditorios pequeños, en estudios o en apartamentos, un sonido "grande" causa molestia al oído cultivado, sin importar la habilidad en el arco. Estas conclusiones se basan en la comprensión general del significado de la palabra "grande" aplicada al sonido del violín. Así, empleado, un sonido "grande" significa un volumen marcado de sonido combinado con una intensidad marcada de sonido. Esta combinación es rara; y, según mi mejor observación, es rara debido a la rareza de la madera de la caja de resonancia que posee una acción de resorte superlativa. Esta afirmación se basa en repetidos fracasos para producir un sonido "grande" a voluntad. Por lo tanto, con todos los demás factores en su mejor momento para aumentar el volumen y la intensidad del sonido, sin embargo, la falta de una acción de resorte superlativa
236

La acción elástica de la madera de la caja de resonancia puede anular un sonido "potente".
En situaciones donde el espacio es proporcional, un sonido de violín "grande" puede resultar agradable. Es un hecho que el ruido es desagradable cuanto más cerca está. También es un hecho que la onda sonora, en todos los violines, viaja solo a distancias relativamente cortas, excepto en casos de timbre amaderado. Dado que los violines con un timbre marcadamente amaderado carecen de valor, no los incluyo en mi afirmación de que el violinista solista no debe temer presentarse en auditorios grandes con un violín moderno de sonido "grande", incluso cuando dicho sonido sea algo ruidoso a corta distancia. Existen algunos violines modernos que combinan un timbre "rico" con un volumen e intensidad notables. Al escucharlos a distancia, he observado que pocos oídos son lo suficientemente agudos como para distinguir alguna diferencia entre el timbre de estos violines modernos y el de violines antiguos que, sin embargo, poseen potencia.
El sonido del violín siempre ha estado, y sin duda seguirá estando, separado en dos clases debido a dos factores irreconciliables, como:
El factor empresarial.
El factor estético.
La razón de ser del factor comercial no tiene otro objetivo que el beneficio; y sus ingresos se calculan en centésimas de dólar.
El factor estético existe únicamente para el placer humano que produce el sonido bello; su contribución es una donación, jamás contabilizada. La diferencia entre estos dos factores es tan grande como la que existe entre Oriente y Occidente.
237

El factor comercial se remonta a tiempos anteriores a Da Salo. El factor estético se remonta a Cremona. El factor comercial siempre ha existido, y sin embargo es numéricamente más fuerte. El factor comercial exige el mayor volumen posible y la mayor intensidad posible del sonido del violín. El factor estético exige la mayor belleza posible en el sonido del violín; sin embargo, para satisfacer la estética
El factor ético es un asunto difícil. El ismo no calcula el costo.	Estética del violín
Para los mortales comunes, el esteticismo del violín a veces parece rozar la locura. La potencia del esteticismo del violín solo es superada por la potencia de la religión. La potencia del esteticismo del violín dio origen al Thomas Hall. No hace falta decir que dicho Hall se encuentra en Chicago. El mundo se enteró de este hecho por la corriente eléctrica.
Ni en Nueva York, ni en Boston, ni en Londres, ni en París, ni en Berlín, ni en Viena, ni en Milán, ni en todo el mundo existe un ejemplo comparable de la potencia del esteticismo del violín antiguo.
Solo el poder de la religión puede rescatar a un millón similar, en parte de debajo de los pies de mujeres humildes; y tales mujeres solo existen en Chicago.
Tras permanecer muchos años en el Auditorio, la Orquesta Thomas no encontró motivo de queja, salvo el hecho de que gran parte de su público no podía oír el sonido del primer violín. Tal queja estaba bien fundada. En los asientos más alejados, la notable ausencia del sonido del primer violín era motivo suficiente de decepción; y mi corazón clamaba contra el corazón...
238

la falta mostrada al forzar estos viejos violines a una situación que despierta desprecio por su debilitamiento
potencia tonal. Hay tres remedios para evitar tal decepción, a saber: 1. Reemplazar la alfombra amortiguadora de fibra de madera, polvo de madera o suciedad dentro de estos violines antiguos con una superficie permanente y perfectamente lisa.
2. Mediante la sustitución por violines modernos con la potencia tonal adecuada.
3. Reduciendo la capacidad de asientos del auditorio.
Con tan solo un aumento del 35 por ciento en la intensidad del sonido del primer violín, la orquesta de Thomas podría haberse quedado en casa en el Auditorio por un período indefinido, evitando así las aceptaciones de los músicos de relleno.
Naturalmente, la sustitución por violines modernos con la potencia tonal adecuada se presenta como el primer remedio para la decepción de no escuchar el sonido del primer violín. Aquí radica toda la dificultad en este caso. Esta sugerencia apunta al uso del violín moderno para la interpretación de partituras de Haydn, Mozart y Beethoven, una posibilidad que el Dr. Thomas no admite. Sobre su firma, el Dr. Thomas afirma que "los mejores violines de Cremona, junto con el arco Tourte, inspiraron las obras maestras de Haydn y Mozart". A esta afirmación no hay disenso; pero a su afirmación de que el violín de Cremona sigue siendo un vehículo necesario para la interpretación de partituras maestras, sí hay disenso; y tal disenso proviene de todo aquel que se siente decepcionado.
239

mecenas. Si esos compositores mayores hubieran estado presentes en el ensayo en el Auditorio, sin duda también se habrían sumado a esa disidencia.
Quizás no haya nada más establecido sobre los "mejores" violines de Cremona (el Stradivarius, al que el Dr. Thomas alude particularmente) que el hecho de que los "mejores" se desgastan por el calor y la humedad, y que los "siguientes mejores" están tan deteriorados que apenas tienen valor para interpretar composiciones, ya sean grandes o pequeñas. Es seguro asumir que ninguno de estos violines antiguos posee el mismo vigor tonal que cuando inspiraron a los Haydn, Mozart y Beethoven de hace cien años.
Se optó por el tercer remedio para superar la decepción. Este ejemplo de devoción a una ilusión costó un millón; pero ese millón no costó más que agacharse para recogerlo. Según toda la lógica histórica y filosófica, dentro de quince años surgirá otra necesidad de reducir la capacidad de los recintos.
En vista de la abundancia de violines modernos que poseen un tono "rico" combinado con una amplia potencia sonora, la afirmación del Dr. Thomas sobre los violines de 200 años como "vehículos necesarios" me parece un esteticismo violinístico llevado al extremo.
Esta digresión se hace con un propósito. Primero: para dejar claro el error de afirmar que el violín de 200 años es un vehículo "necesario" para la interpretación de composiciones de 100 años. Segundo: para dejar claro que el oyente está en mejor posición.
240

ción para juzgar los valores tonales del primer violín. Basándonos en las conclusiones de la declaración del Dr. Thomas, nos vemos obligados a suponer que la interpretación de las partituras de Haydn, Mozart y Beethoven debe cesar con "los mejores violines de Cremona". Estoy convencido de que esta postura del Dr. Thomas es el resultado de un esteticismo del violín antiguo llevado al extremo; y, además, que su postura se demostrará insostenible en un futuro próximo. Que la declaración del Dr. Thomas excite desprecio por el violín moderno es claramente evidente; también, que su insistencia en emplear solo violines antiguos por la orquesta de Thomas provocó desprecio por "los mejores Cremona", como lo demuestran las decepcionantes experiencias de un ejército de mecenas. No cabe duda de que las interpretaciones de la gran orquesta están destinadas al disfrute de todo el público, y no al disfrute especial del director y los ocupantes de las primeras filas. Tampoco cabe duda de que el oyente se sentirá decepcionado al no escuchar el sonido del primer violín en el conjunto. Hablo solo por mí cuando digo que la ausencia del sonido del primer violín en una orquesta es una interpretación poco artística, ya sea que la partitura tenga 100 años o apenas una hora. Al decir esto, no quiero decir que se pueda permitir un desequilibrio excesivo en las cuerdas. El mero volumen del sonido del violín sin dulzura resulta ofensivo tanto en violines solistas como en orquestales. Lo que considero las características más valiosas en el sonido del violín es un volumen moderado combinado con dulzura, uniformidad en la potencia y una marcada
241

Intensidad. Esta combinación de cualidades tonales no puede resultar ofensiva para el oído; además, existe una amplia oferta de violines modernos con esta combinación de valores tonales para las necesidades de las grandes orquestas que interpretan obras de Haydn, Mozart y Beethoven, o de cualquier otro compositor de renombre. Esto nos lleva a los detalles para lograr la máxima potencia sonora del violín.
Según mi observación, cada uno de los siguientes factores debe estar en su punto óptimo para lograr la máxima potencia tonal del violín:
1. Madera de la caja de resonancia: Calidad de muelles superlativa.
2. Graduación: Forma que asegura la mayor fuerza
del soplo sobre el aire contenido.
3. identidad.	Atrás:	Densidad, textura fina, estructura suficiente
4. Forma arqueada: Forma que produce la máxima concentración del movimiento de las ondas sonoras en las salidas.
5. Barra: Densidad media, modelada para la mayor
acción de resorte, posición.
6. Salidas: Área, posición.
7. Superficies interiores: Lisura permanente.
8. Costillas: Profundidad, rigidez.
9. Bloques y revestimientos: Solidez.
10. Poste: Masa, densidad, longitud, ajuste, posición.
11. Capacidad cúbica del cuerpo.
12. Barniz: Cantidad, calidad. 13. Puente: Masa, densidad, altura, extensión, ornamentación, posición. 14. Diapasón: Largo, ancho, altura, obli- 242

paz de la superficie inferior a la superficie exterior de la caja de resonancia.
15. Cuerdas: Diámetro, proporción, número de
hebras, torsión, material.
16. Vara del arco: Longitud, peso, equilibrio, curvatura, acción de resorte.
17. Pelo de arco: Equinos machos, uniformidad, método de
ensamblaje. 18. Resina: Calidad, cantidad.
19. Condiciones meteóricas: Temperatura moderada,
humedad normal.
20. Elevación sobre el nivel del mar.
21. Brazo de arco: longitud, condición, entrenamiento, entusiasmo.
De esta lista de factores, coloco la madera de la caja de resonancia en primer lugar; y esta posición destacada se debe enteramente a que no he podido producir la máxima potencia tonal del violín sin una madera de caja de resonancia con una elasticidad excepcional, y también a que no he podido encontrar dicha madera todos los meses, ni todos los años, ni todas las décadas. De hecho, solo unas pocas veces en cincuenta años de búsqueda he encontrado dicha madera; mientras que encontrar madera dura con densidad y textura fina ha sido comparativamente muy fácil. Basta con una observación limitada para determinar que ni la madera más blanda ni la más dura producen la máxima resonancia; pero una observación más amplia no puede predeterminar que una muestra de madera determinada producirá la máxima resonancia en el violín terminado. Solo después de la finalización y las pruebas se puede conocer el grado de resonancia de cualquier violín.
243
t

Independientemente del fabricante, dos muestras de madera para la caja de resonancia pueden parecer idénticas; sin embargo, en el violín terminado, puede haber una gran diferencia en la resonancia, incluso cuando los detalles de construcción son similares. Por lo tanto, lograr la máxima potencia sonora en el violín se convierte en una tarea difícil, y dicha dificultad se debe a las variaciones inherentes a la elasticidad de la madera.
Resulta interesante observar que los violines de máxima resonancia y con un sonido rico son una rareza, incluso después de 400 años de esfuerzos por producirlos. Evidentemente, existe algún obstáculo que impide su fabricación; de lo contrario, habría una gran cantidad de violines disponibles. La causa de esta rareza puede no ser tan evidente para otros como lo es para mí; pero todos coinciden en que estos violines alcanzan precios elevados. Dado que hoy son escasos, podemos suponer con seguridad que lo seguirán siendo mañana y durante un período indefinido. Es evidente que la causa de esta escasez escapa al control humano.
Es un hecho de demostración cotidiana que el temple de los resortes, ya sea en madera o metal, es algo que se conoce solo como resultado de la acción después de la prueba; y, dicho resultado no puede predeterminarse. Aunque el templador de resortes de acero puede producir grados de temple "altos" o "bajos" a voluntad, producir cualquier número dado de resortes de acero que posean precisamente una acción similar es una tarea difícil.
244

Se trata de una cuestión compleja, cuya dificultad se debe a dos causas: (1) Variaciones en la disposición molecular de diferentes muestras de acero; (2) Incapacidad para predeterminar dichas variaciones. En la práctica, se demuestra claramente que variaciones imperceptibles en la densidad o dureza del resorte de acero producen variaciones perceptibles en su acción. Asimismo, el valor del resorte depende de su temple, dureza o densidad, según el término que se elija.
Naturalmente, nos preguntamos cuál sería el resultado si a la madera de la caja de resonancia del violín se le diera "temperamento" hasta el límite del deseo. ¡Sombra de Coloso! Entonces, ¿por qué el violín actual de tono "grande" no serviría para el bebé? Y, el auditorio más grande no serviría para un escenario; y, el género de
La música tendría que cambiarse a "él"; "mataría a "ella"; y "él" tendría que ser sordo. ¿Encontrar usuarios?
Créelo. Tal como veo este asunto, la Naturaleza sabiamente limita al hombre
posibilidades en el violín.
Lo que diré acerca de las apariencias de la madera de la caja de resonancia que prometen máxima resonancia de tono, combinada con una calidad de tono "rica", se presenta solo como la observación y conclusión de un individuo; y, de ninguna manera dicha conclusión se presenta como una regla estricta. En cuanto al tema de la resonancia en la madera de la caja de resonancia, debo confesar que infinitas sorpresas han aguardado a mi arco. Se observa que la máxima resonancia de
245

El timbre se combina con la riqueza tonal. Deseo que esta combinación se mantenga en primer plano, pues la fuerza sin dulzura reduce el valor tonal del violín a la insignificancia. En mi opinión, ningún hecho sobre el timbre del violín es más evidente que el hecho de que la veta de la caja de resonancia puede ser tan densa que imposibilite un timbre rico, mientras que la mera fuerza tonal se consigue con facilidad.
De las maderas de pino más blandas se puede obtener un tono rico con bastante certeza; pero, de esas mismas maderas, es imposible lograr una potencia tonal marcada. Cabe destacar que un gran volumen de sonido no implica una gran intensidad tonal. Se puede obtener un gran volumen de sonido de la madera más blanda, pero nunca una gran intensidad tonal.
La enorme diferencia entre volumen e intensidad del sonido queda patente en la prueba a larga distancia al aire libre. Esta prueba demuestra de forma concluyente que un gran volumen no es necesario para lograr una gran potencia sonora en el violín. De hecho, no he observado que un gran volumen y una gran intensidad de sonido coexistan en el mismo violín.
En mi opinión, las características físicas de la madera de la caja de resonancia que producen la máxima resonancia combinada con una calidad de tono "rica" son:
Por lo tanto: Género pino.
Grado de madurez del árbol antes de ser talado.
Muestra libre de duramen y albura, nudos, curvaturas, grietas y manchas descoloridas.
Veta perfectamente recta, uniforme en ancho, ni
246

ni muy ancho ni extremadamente estrecho. Color rojo amarillento.
[*]Tejido conectivo, entre partes densas del grano, bastante blando y flexible.
La parte dura de la veta está claramente marcada. La madera se agrieta con facilidad. Las grietas siguen una línea recta.
Las virutas son más bien quebradizas que resistentes. Se deshilachan fácilmente. Resulta difícil conseguir una superficie lisa. La transmisión de luz artificial es media. La placa es más bien flexible que rígida.
La placa doblada vuelve rápidamente a su punto de reposo.
Sin protección, el calor y la humedad desarrollan rápidamente una capa de fibras de madera.
Desintegración relativamente rápida. Densidad media del grano.
Así son las cualidades físicas de la madera ideal para mi caja de resonancia. Dicha madera, con los modelos correctos de curvatura, la graduación adecuada, el área y la posición correctas de las aberturas, y con superficies interiores permanentes y perfectamente lisas, proporciona una amplia intensidad de tono para todos los usos prácticos del violín, además de producir un tono "rico". Como recordará, el tono "rico" del violín depende de la presencia de armónicos superiores y armónicos bajos, o tonos resultantes, y de que dichos armónicos no pueden aparecer en la caja de resonancia de madera densa.
grano.
Hay grados de densidad entre la densidad media y la extremadamente densa que poseen consi-
247

valor verificable; y, en situaciones donde el ruido
abunda, la caja de resonancia de más de medio
La densidad posee el mayor valor. En tales situaciones, esos tonos etéreos llamados armónicos se aniquilan; incluso el tono principal puede tener que luchar por existir; por lo tanto, la caja de resonancia que posee la mayor potencia en acción de resorte da la mayor satisfacción en presencia de
ruido.
Pero no encuentro que un aumento en la densidad de la madera de la caja de resonancia se traduzca necesariamente en un aumento de la potencia del sonido. Por el contrario, observo una disminución de la potencia del sonido con densidades extremas. Este resultado podría ser una suposición; y, dado que un aumento de la densidad conlleva un aumento de la rigidez, y un aumento de la rigidez conlleva una disminución de la amplitud de oscilación, y una disminución de la amplitud de oscilación ejerce una menor fuerza sobre el aire contenido, se produce un debilitamiento del sonido.
El método de graduación tiene mucho que ver con la resonancia del violín. En este momento, la graduación de la caja de resonancia es relevante. El grosor de la caja de resonancia puede ser tan grande o tan pequeño que cause debilidad en el tono. Dado que las cifras precisas para el grosor de la tapa del violín no son guías fiables, no se presentan dichas cifras. En mi experiencia, se demuestra claramente que el grosor de la caja de resonancia, que produce la mayor fuerza de los golpes sobre el aire contenido, varía con cada muestra.
de caja de resonancia wocd; y, por esta razón, las cifras fijas para los espesores son guías poco fiables que con frecuencia conducen al desastre 248

que al éxito. Es evidente que la rigidez de la caja de resonancia debe determinarse por la fuerza de las cuerdas, ya que los golpes de las cuerdas son el único medio para activar la caja de resonancia. Por lo tanto, primero debe determinarse el calibre de las cuerdas y, posteriormente, la rigidez de la caja de resonancia de cada violín debe reducirse al grado que permita aplicar el máximo golpe al aire contenido; y dicho grado de rigidez solo puede determinarse mediante ensayo y error. Evidentemente, para tener certeza en este asunto, no hay alternativa.
La clave está en saber cuándo cada caja de resonancia está ejerciendo su máxima potencia. En este sentido, no conozco ninguna regla fija, salvo la que impone la experiencia. Puedo afirmar que una amplia experiencia en la reducción de la rigidez de las cajas de resonancia lleva a pecar de precavido en lugar de arriesgarse a un desastre por una reducción excesiva. Según mi experiencia, una vez alcanzado el grado exacto de rigidez, la eliminación de virutas aún más finas debilita la fuerza del golpe sobre el aire contenido.
Para asegurar el máximo impacto sobre el aire contenido, es necesario considerar con sumo cuidado la magnitud de la acción vibratoria, tanto normal como transversal, en la caja de resonancia.
(Véase el Apéndice para la explicación de la vibración normal y transversal).
Como estas acciones se producen por los golpes de las cuerdas, y como las cuerdas varían en diámetro y peso, por lo tanto, la rigidez de la caja de resonancia también varía.
249

La fuerza de la caja de resonancia debe variar según la intensidad de las cuerdas para lograr la máxima potencia tonal. Dado que la fuerza en la caja de resonancia es menor que la de la cuerda A, la de la A menor que la de la de la de D, y la de la de la de D menor que la de la de G, no es razonable esperar que la máxima vibración se logre mediante un método de gradación de la caja de resonancia que proporcione la misma rigidez de la madera a todas las cuerdas.
Al hablar de la acción vibratoria que se propaga a través de la veta de la caja de resonancia, los términos "transversal" y "tangencial" pueden emplearse indistintamente, ya que, en la práctica, ambos tienen el mismo significado. En la práctica, aumentar la distancia recorrida por la vibración transversal en la caja de resonancia del violín incrementa el volumen del sonido; y dicho aumento puede llegar a un punto en el que se produzca una pérdida de intensidad. Este punto es de gran importancia para lograr la máxima potencia sonora, ya que un gran volumen no implica una gran potencia sonora. La gran potencia sonora se refiere a la distancia recorrida por el sonido. En la caja de resonancia del violín, la distancia recorrida por la vibración transversal depende de tres factores.
factores:
1. El poder de las cuerdas.
2. Rigidez de la placa a lo largo de líneas perpendiculares a la junta central.
3. Densidad del grano.
Evidentemente, es necesario determinar primero el tamaño o calibre de las cuerdas. Segundo, disminuir la rigidez a lo largo de líneas perpendiculares al centro, uniéndolas a ese grado, permitiendo que las cuerdas impulsen la vibración transversal lo más lejos posible del centro.
250

unir como se desee. Algunos fabricantes de violines reducen la rigidez en este sentido hasta un grado que permite que la fuerza de las cuerdas impulse la vibración transversal hasta el borde de la placa. Este plan invariablemente aumenta el volumen del tono; pero también disminuye la intensidad del tono; y la razón de tal pérdida de intensidad se debe claramente a la dispersión de la fuerza de las cuerdas. He observado que, para producir la máxima potencia tonal del violín, la fuerza de las cuerdas debe concentrarse en lugar de dispersarse; y por lo tanto, limito la distancia recorrida por la vibración transversal. [Es necesario tener en cuenta que empleo cuerdas de calibre 2, porque, si se emplearan cuerdas más gruesas, habría mayor fuerza;
Por lo tanto, podría permitirse de forma segura una mayor distancia de recorrido para la vibración transversal.]
Existen dos métodos para determinar la distancia que debe recorrer la vibración transversal en la caja de resonancia del violín una vez que su grosor se ha reducido casi al grado final, de la siguiente manera:
1. Eliminando la madera de la superficie interior antes del montaje.
2. Mediante la eliminación de la madera de la superficie exterior después del montaje y el tendido.
Con el primer método, los resultados dependen de conjeturas al trabajar con madera nueva y sin probar. Con el segundo método, los resultados son seguros, ya que, con las cuerdas afinadas, se puede comprobar su fuerza hasta que la vibración transversal alcance cualquier punto deseado en líneas perpendiculares a la unión central.
Para concentrar la fuerza de la cuerda sobre el sonido-
251

En mi experiencia, el siguiente método para disminuir el grosor ha sido el más exitoso; porque, con este método, pude limitar la distancia recorrida por la vibración transversal en el punto que produce lo que yo llamo un volumen de tono suficiente sin la mayor dispersión de la fuerza de la cuerda; y el resultado ha sido un volumen de tono moderado con una marcada intensidad de tono, una combinación que invariablemente alcanza la mayor distancia en la prueba de larga distancia al aire libre.
Así pues: El mayor espesor se alcanza en la posición del puente; desde allí, el espesor disminuye gradualmente hasta un punto situado a mitad de camino entre la posición del puente y los extremos de la placa; desde dicho punto, el espesor aumenta gradualmente hasta los extremos de la placa; y, desde dicho punto intermedio, en líneas perpendiculares a la junta central, el mismo espesor en dicho punto se mantiene a lo largo de dichas líneas laterales hasta un punto situado a una pulgada y un cuarto a la derecha, y hasta un punto situado a una pulgada y un cuarto a la izquierda de la junta central; y, desde estos puntos, los espesores también aumentan gradualmente en ambas direcciones desde dichas líneas laterales.
A continuación, una vez ensamblado el violín y afinadas las cuerdas, se comprueba la fuerza de estas. Intencionadamente, la caja de resonancia queda demasiado rígida al ensamblarla, como medida de precaución.
Ahora llega el momento de un trabajo en el que la ciencia no puede sustituir al sentido musical; ni la riqueza mundial puede comprar ni una pizca de sentido musical. El sentido musical no es en absoluto una mercancía. El sentido musical no se puede tomar prestado.
Ahí está el problema.
252

Puedo tomar prestados los libros de mi amigo y olvidarme rápidamente de devolverlos; pero nunca tendré la oportunidad de olvidar la devolución de su sentido musical. Al probar el tono del violín recién ensamblado, y encontrar que la duración del tono es demasiado corta para mi gusto, entonces aumento la distancia recorrida por la vibración normal; y, logro tal resultado disminuyendo el grosor de la caja de resonancia desde la línea lateral media indicada hasta los extremos de la placa; y, si encuentro que el volumen del tono aún es demasiado pequeño, entonces aumento la distancia recorrida por la vibración transversal; y logro este resultado disminuyendo el grosor en, por encima, por debajo y fuera de los puntos laterales indicados; y, para dar una oportunidad igual a la fuerza en las cuerdas A y E, se disminuye el grosor inmediatamente debajo de esas cuerdas.
Al trabajar con madera nueva y sin probar, reitero la importancia de ser precavido en lugar de disminuir el grosor hasta que se produzca una pérdida de tono. Esta precaución es necesaria debido al aumento de flexibilidad del tejido conectivo, inevitable tras el uso. He visto cómo este aumento de flexibilidad puede arruinar los valores tonales en tan solo dos años de uso. Al regular el grosor de la caja de resonancia con el objetivo de obtener la máxima potencia tonal, es importante tener en cuenta que la acción simpática entre los cuerpos vibrantes disminuye con el cuadrado de la distancia; por lo tanto, el grosor en la mitad inferior de la caja de resonancia debe ser menor que en la mitad superior. En cuanto al trabajo sobre la propia caja de resonancia,

En lo que respecta a estos detalles, se incluyen las características principales; pero, como muestra la lista anterior, hay muchos otros factores necesarios para la producción de la máxima potencia tonal del violín. Dado que los factores más importantes de la lista se han presentado en páginas anteriores, no es necesario extenderse en este momento, salvo decir que, en mi experiencia, ningún factor de la lista puede descuidarse sin provocar un desastre. Hay numerosos ejemplos, además del violín, que muestran los efectos desastrosos debidos a la dispersión de la fuerza. La trompa, demasiado ligera en metal, cede, tiembla, "dobla" ante la fuerza de la
Los pulmones del músico pierden potencia tonal. En lugar de controlar la energía elástica de las moléculas de aire en su interior, se ven controlados por sí mismos.
Una vez más, el silbato de vapor solo produce un tono armónico cuando la presión del vapor supera con creces la rigidez del silbato.
Una vez más, el piano, con una caja de resonancia demasiado ligera, produce un sonido débil; un resultado que se debe claramente a la dispersión de la fuerza.
Una vez más, la pelota, lanzada contra la lona que cuelga, cae directamente al suelo; un resultado debido a la dispersión de la fuerza.
Nuevamente, la voz fuerte, dirigida hacia el diafragma del teléfono, provoca que el receptor la escuche sin claridad; un resultado debido a la dispersión de la fuerza.
De nuevo, al orador, al aire libre, solo se le puede oír a distancias limitadas debido a la dispersión de la fuerza.
Nuevamente, el arma de caza, al retroceder, pierde alcance. De estos ejemplos, se puede decir que la fuerza es demasiado
254

excelente.
Cierto desde un punto de vista, falso desde otro. El violinista, que siempre busca una mayor potencia sonora que la que produce el peso de su arco, no se siente a gusto con un violín que requiere una presión vigorosa. Para complacer a dicho violinista, el luthier debe fabricar un violín ligero; y este violín no resulta del agrado del violinista que posee vigor y entusiasmo en el arco, porque, al ser sometido a una presión vigorosa, vibra hasta tal punto que permite una dispersión excesiva de la fuerza en las cuerdas, lo que resulta en un sonido débil.
La pérdida de potencia no es el único problema que sufre el violín por el uso de madera demasiado ligera. Un problema aún más desastroso es la calidad del sonido amaderado; el tipo de sonido que se obtiene con madera dura cuando su rigidez es excesiva.
reducido.
He visto violines, construidos con las mejores intenciones, arruinados por una excesiva reducción de la rigidez de la tapa de madera dura. Estos violines, con una ligera presión del arco y cuerdas finas, pueden producir un sonido aceptable; pero, con una presión vigorosa del arco, el sonido es insoportablemente amaderado. Como instrumento musical, un violín así cae en la insignificancia.
